\section{Cusps, McShane gaps, and the orthospectrum}
\label{sec:cusps}

\subsection{Which geodesics enter a cusp?}

Throughout, the surface is oriented, hyperbolic of curvature $-1$, and
has no cone points. The statements that simple closed geodesics avoid a sufficiently deep cusp
and that simple cusp-ended geodesics enter the cusp are compatible. The
first concerns a returning excursion; the second concerns an actual end.

Here is an elementary version of the distinction. Normalize a cusp
stabilizer in the upper half-plane to $z\mapsto z+1$. A geodesic with two
finite ideal endpoints $u<v$ is a semicircle of radius $(v-u)/2$. If
$v-u>1$, its endpoints interlace those of its translate:
\[
 u<u+1<v<v+1.
\]
The two geodesics intersect, so the projected geodesic is not simple.
Therefore a simple geodesic with two finite endpoints cannot reach height
greater than $1/2$. A geodesic reaching farther into this lifted cusp
while remaining simple must have $\infty$ as an endpoint, and is vertical.
It is perpendicular to every horocycle based at that endpoint.

Choose $Y>1/2$ large enough that the horocycle $H=\{y=Y\}/\langle z\mapsto
z+1\rangle$ bounds an embedded sufficiently deep cusp. Simple closed
geodesics never cross $H$. Complete simple geodesics with a cusp end do
cross it, normally, and continue indefinitely toward that end.

\subsection{The family used in the gap argument}

For this gap discussion assume that the surface is complete, has finite
area and at least one cusp, and has no geodesic boundary. A surface with
geodesic boundary has additional families of cusp-issued rays that exit
through it, and needs the corresponding boundary gap terms.

At a point $x\in H$, take the normal geodesic directed away from the
cusp. A lift is
\begin{equation}
 \widetilde\gamma_x(t)=x+iYe^{-t}.
 \label{eq:cuspray}
\end{equation}
Its backward end tends to the cusp. Its forward continuation enters the
rest of the surface. The relevant question is whether that whole
continuation remains simple; the short portion inside the cusp is
automatically simple.

McShane's gap argument partitions the horocycle using the behavior of
these cusp-issued geodesics. Almost every direction eventually develops
the relevant intersection pattern. The exceptional simple directions
have measure zero and include cusp-to-cusp arcs and rays spiraling toward
closed geodesics or minimal laminations
\cite{McShane}\cite[Theorem 3, p.~4, and the paragraph after
Corollary 5, p.~5]{McShanePants}.
In some conventions a gap contains an isolated simple cusp-to-cusp
direction; its interior is therefore not literally made entirely of
nonsimple directions. The measure statement is unaffected.

For a cusp pants contribution whose two other cuffs are interior, with
lengths $a,b$ (allowing their identification on the torus), the combined
normalized gap is
\begin{equation}
 \frac{2}{1+e^{(a+b)/2}}.
 \label{eq:cuspgap}
\end{equation}
For example, this is the limit as $L\downarrow0$ of the geodesic-boundary
combined gap divided by $L$:
\[
 \frac{2}{L}\log
 \frac{e^{L/2}+e^{(a+b)/2}}{e^{-L/2}+e^{(a+b)/2}}
 \longrightarrow \frac{2}{1+e^{(a+b)/2}}.
\]
The boundary gap is Mirzakhani's $\mathcal D$ contribution
\cite[Lemma 3.1, equation (3.1), and Theorem 4.2]{MirzVolumes}. This calculation fixes the
normalization; the null-set and pants-indexing arguments are separate
external inputs.
On the punctured torus the two cuffs are copies of one simple curve
$\alpha$, so $a=b=\ell_\alpha$. The gap contributions sum to one,
which is precisely~\eqref{eq:mcshane}. The simple closed curve occurring
in the formula labels the pants and their gap. It is not the geodesic
crossing the deep horocycle.

\subsection{An exact truncated cusp-arc formula}

Fix a sufficiently small embedded cusp horocycle of length $H>0$; here $H$ now
denotes its length rather than the horocycle as a set. In pants with other
cuffs $a,b$, let $\eta$ be the simple cusp-to-cusp geodesic arc separating
those cuffs. Truncate its two ends at the horocycle and let $d_\eta$ be
the resulting length. With the convention $\lambda_\eta=e^{d_\eta/2}$,
\begin{equation}
 \boxed{e^{d_\eta/2}=\frac{t_a+t_b}{H},
        \qquad t_a=2\cosh(a/2),\quad t_b=2\cosh(b/2).}
 \label{eq:cusplambda}
\end{equation}
The other cuffs may be geodesic or cuspidal, so $a,b\geq0$.
The formula also holds for signed truncated lengths of decorations before
the horocycle has been made small enough; this does not assert that an
arbitrarily large decorated horocycle is embedded.

\begin{proof}
Normalize the cusp parabolic and horocycle to
\[
 P=\begin{pmatrix}1&q\\0&1\end{pmatrix},\quad q>0,\qquad y=Y,
 \qquad H=q/Y.
\]
Choose cuff holonomy lifts $A,B$ with $AB=P$, $\tr A=t_a$ and
$\tr B=-t_b$, and write $c_A$ for the lower-left entry of $A$. Then
\[
 \tr(A^{-1}P)=t_a-qc_A=-t_b,
 \qquad qc_A=t_a+t_b.
\]
The image under $A$ of the horizontal horocycle has Euclidean diameter
$1/(c_A^2Y)$ and is based at $A(\infty)$. The vertical ideal geodesic
from $\infty$ to this point represents the separating return arc.
The upper intersection with that image horocycle is at height
$1/(c_A^2Y)$, so its signed truncated distance is
$\log(c_A^2Y^2)$ and satisfies
$e^{d_\eta/2}=|c_A|Y$. Substituting $H=q/Y$ proves the formula.
\end{proof}

On the once-punctured torus, essential unoriented simple cusp-to-cusp arc
classes correspond to essential unoriented simple closed-curve classes:
each is the unique such class disjoint from the other
\cite[Proposition 14.1(i),(iv)]{Springborn}. Cutting along the
closed curve gives $(a,a,0)$ pants, hence
\begin{equation}
 e^{d_\alpha/2}=\frac{2t_\alpha}{H}
              =\frac{4\cosh(\ell_\alpha/2)}{H}.
 \label{eq:toruscusplambda}
\end{equation}
As a check, each decorated ideal triangle contributes horocyclic corner
lengths $u/(vw)$, $v/(uw)$ and $w/(uv)$ in the convention
$\lambda=e^{d/2}$~\cite[Section 11, equation (30)]{Springborn}.
The two triangles of a torus therefore give
$H=2(u^2+v^2+w^2)/(uvw)$, as in that paper's Section 12.
By~\eqref{eq:toruscusplambda}, the positive trace of the curve disjoint
from edge $u$ is $Hu/2$. Thus the trace conversion is derived here from
the cusp-parabolic calculation, with no additional factor of $\sqrt2$
in the lambda convention.

Substitute $t_\alpha=(H/2)e^{d_\alpha/2}$ into
$\sum_\alpha h(t_\alpha)=1$. This gives a literal orthospectral form of
McShane's identity:
\begin{equation}
 \boxed{
 \sum_{\eta\text{ simple, unoriented cusp-to-cusp}}
 \left(1-\sqrt{1-\frac{16e^{-d_\eta}}{H^2}}\right)=1.}
 \label{eq:cusporthoidentity}
\end{equation}
All square roots use the nonnegative real branch. Indeed,
$16e^{-d_\eta}/H^2=4/t_\alpha^2<1$, so the summands are real and
positive, even when a chosen signed truncated length is negative.
The index set retains the topological arc labels. Scaling the horocycle
length by $k$ changes each truncated length to $d_\eta-2\log k$;
therefore $e^{-d_\eta}/H^2$ and every summand are decoration-independent.

\subsection{The precise relation to Basmajian's identity}

For a compact hyperbolic surface with nonempty geodesic boundary,
Basmajian's boundary-length identity is, with oriented orthogeodesics,
\begin{equation}
 \ell(\partial X)=
 \sum_{\eta\text{ oriented}}2\log\coth\left(\frac{\ell_\eta}{2}\right).
 \label{eq:basmajian}
\end{equation}
The oriented convention and this normalization are also explicit in
\cite[Section 4.2, p.~26]{Parlier}.
If each unoriented orthogeodesic is counted only once, the coefficient is
$4$ instead of $2$. This spectrum includes nonsimple arcs. The gap comes
from the shadow on a boundary lift of another boundary lift, and
hyperbolic trigonometry expresses its length in terms of their common
perpendicular. To see the coefficient directly, take one boundary lift
to be $x=0$ and the other the semicircle with center $\coth d$ and
radius $\operatorname{csch}d$, where $d=\ell_\eta>0$. Its ideal endpoints
are $\tanh(d/2)$ and $\coth(d/2)$. Their perpendicular projections to
$x=0$ have these heights, and the shadow length is
$\log(\coth(d/2)/\tanh(d/2))=2\log\coth(d/2)$.
The opposite orientation records the shadow at the other end. For a
large $d$, the equivalent expression
$2\log(1+2e^{-d}/(1-e^{-d}))$ retains the positive gap without subtracting
nearby numbers.

Equation~\eqref{eq:cusporthoidentity} uses the \emph{simple, truncated
cuspidal orthospectrum}. Generic normal rays in the McShane gaps are not
the canonical returning orthogeodesics, and adding all nonsimple cusp arcs
would change the identity. Thus the orthospectrum interpretation is exact
after the index set, truncation, and weights have been specified.

Parlier's coherent-marking framework gives a systematic bridge between
these boundary and cusp identities~\cite{Parlier}. The common mechanism
is a measure decomposition of boundary initial data, with orthogeodesic
geometry computing gaps. What differs is the chosen stopping or marking
rule and the resulting index set. This is the useful sense in which the
pants, McShane, and Basmajian pictures fit together. A finite selection
of cusp arcs or orthogeodesics is not a verified approximation to these
whole-surface identities unless its omitted classes are also controlled.

\paragraph{Computational support.}
Groups G9--G10 of \path{experiments/quantum_geodesics/README.md}
check exact parabolic/decoration and ideal-triangle identities,
alongside finite local shadow and orientation-factor comparisons.
Signed truncation and stable small positive shadow formulas are
retained. These are not a discrete-holonomy certificate, a complete
orthospectrum enumeration or a numerical proof of its global identity.
The result path is \path{artifacts/quantum_geodesics/result.json},
with claim pointers in \path{book/claims/sections17-22.json}.
